\documentclass[12pt,letterpaper]{article}

\usepackage[spanish,es-tabla]{babel}
\usepackage[utf8]{inputenc}
\usepackage[T1]{fontenc}
\usepackage{geometry}
\usepackage{graphicx}
\usepackage{float}
\usepackage{longtable}
\usepackage{array}
\usepackage{booktabs}
\usepackage{hyperref}
\usepackage{xcolor}
\usepackage{listings}
\usepackage{enumitem}
\usepackage{titlesec}
\usepackage{caption}
\usepackage{pdflscape}
\usepackage{tabularx}
\usepackage{xurl}

\geometry{margin=2.4cm}
\hypersetup{
    colorlinks=true,
    linkcolor=blue!55!black,
    urlcolor=blue!55!black,
    citecolor=blue!55!black
}

\definecolor{codebg}{RGB}{248,248,248}
\definecolor{codeframe}{RGB}{220,220,220}
\definecolor{restblue}{RGB}{10,43,120}

\lstdefinelanguage{json}{
    basicstyle=\ttfamily\small,
    showstringspaces=false,
    breaklines=true,
    morestring=[b]",
    stringstyle=\color{blue!50!black},
    literate=
     *{0}{{{\color{black}0}}}{1}
      {1}{{{\color{black}1}}}{1}
      {2}{{{\color{black}2}}}{1}
      {3}{{{\color{black}3}}}{1}
      {4}{{{\color{black}4}}}{1}
      {5}{{{\color{black}5}}}{1}
      {6}{{{\color{black}6}}}{1}
      {7}{{{\color{black}7}}}{1}
      {8}{{{\color{black}8}}}{1}
      {9}{{{\color{black}9}}}{1}
      {:}{{{\color{black}:}}}{1}
      {,}{{{\color{black},}}}{1}
      {\{}{{{\color{black}\{}}}{1}
      {\}}{{{\color{black}\}}}}{1}
      {[}{{{\color{black}[}}}{1}
      {]}{{{\color{black}]}}}{1},
}

\lstdefinelanguage{bash}{
    keywords={docker,compose,up,down,logs,python,pytest,powershell,ruby,bin,rails,true,false,SELECT,FROM,WHERE},
    sensitive=true,
    comment=[l]{\#},
    morestring=[b]",
}

\lstdefinestyle{codestyle}{
    backgroundcolor=\color{codebg},
    frame=single,
    rulecolor=\color{codeframe},
    basicstyle=\ttfamily\small,
    breaklines=true,
    columns=fullflexible,
    showstringspaces=false,
    tabsize=2
}
\lstset{style=codestyle}

\setlist[itemize]{noitemsep,topsep=3pt}
\setlist[enumerate]{noitemsep,topsep=3pt}
\titleformat{\section}{\Large\bfseries\color{restblue}}{\thesection}{0.7em}{}
\titleformat{\subsection}{\large\bfseries}{\thesubsection}{0.7em}{}
\titleformat{\subsubsection}{\normalsize\bfseries}{\thesubsubsection}{0.7em}{}

\urlstyle{tt}

\def\endpoint#1#2{\texttt{#1 #2}}

\def\codefile#1{\nolinkurl{#1}}

\def\endpointcell#1#2{%
  \textbf{#1}\\
  \codefile{#2}%
}

\def\methodget{\textbf{\color{blue!60!black}GET}}
\def\methodpost{\textbf{\color{green!45!black}POST}}
\def\methodput{\textbf{\color{orange!80!black}PUT}}
\def\methoddelete{\textbf{\color{red!70!black}DELETE}}

\title{
    \textbf{Documentación Técnica Completa}\\
    Sistema Distribuido de Restaurantes\\
    \large FastAPI, Keycloak, PostgreSQL, MongoDB Sharding, Redis, Elasticsearch, Docker, Kubernetes y Plataforma Analítica \texttt{dataW}
}
\author{Cesar Gamboa \and Daniel Morera}
\date{Junio 2026}

\begin{document}

\maketitle
\thispagestyle{empty}

\begin{abstract}
Este documento consolida la documentación funcional, técnica y operativa del proyecto \textit{Restaurante-DB-I-Proyecto}. Describe la arquitectura física y lógica, el comportamiento de cada servicio, las tecnologías utilizadas, los contratos de API, los modelos de datos, los flujos de información y la carpeta \codefile{dataW}, que implementa el componente analítico con Hive, Spark, Superset y Airflow. La documentación está basada en los archivos del repositorio, los diagramas y los requerimientos del proyecto.
\end{abstract}

\newpage
\tableofcontents
\newpage

\section{Resumen del Programa}

El proyecto implementa una plataforma distribuida para gestionar restaurantes, usuarios, menús, mesas, pedidos, reservaciones, búsqueda de productos y análisis de datos. El sistema se divide en dos grandes dominios:

\begin{itemize}
    \item \textbf{Dominio transaccional}: API REST en FastAPI, autenticación con Keycloak, persistencia seleccionable entre PostgreSQL y MongoDB shardeado, caché con Redis, búsqueda con Elasticsearch y entrada unificada con Nginx.
    \item \textbf{Dominio analítico}: carpeta \codefile{dataW}, que contiene un Data Warehouse en Hive, procesamiento con Spark, dashboards en Superset y orquestación con Airflow.
\end{itemize}

El sistema corre principalmente sobre Docker Compose. También se incluye una definición de despliegue Kubernetes para los servicios principales. La API puede trabajar contra PostgreSQL o MongoDB según la variable \codefile{DATABASE\_ENGINE}, manteniendo una capa DAO común para desacoplar rutas y persistencia.

\section{Cómo leer esta documentación}

Esta documentación está organizada con un estilo similar a un portal de documentación para desarrolladores, se tomo como referencia la wed oficial de Spotify. Primero se presenta una visión general, luego una guía de inicio, conceptos clave, guías prácticas y finalmente una referencia de API detallada. En la documentación oficial de Spotify, la Web API separa \textit{Overview}, \textit{Getting started}, \textit{Concepts}, \textit{Tutorials}, \textit{How-Tos} y \textit{Reference}; cada endpoint de referencia muestra método, ruta, autenticación, parámetros, respuesta y códigos de error.

\subsection{Mapa de navegación}

\begin{longtable}{p{0.24\textwidth}p{0.66\textwidth}}
\toprule
\textbf{Sección} & \textbf{Qué resuelve}\\
\midrule
\textbf{Overview} & Explica qué es el sistema, qué problema resuelve y cuáles son sus componentes.\\
\textbf{Getting started} & Indica cómo levantar el proyecto, obtener token y hacer la primera llamada.\\
\textbf{Concepts} & Define tokens, roles, motores de base de datos, caché, búsqueda, sharding y pipeline analítico.\\
\textbf{How-Tos} & Muestra tareas comunes: registrar usuario, crear restaurante, reindexar búsqueda, correr Airflow y validar \codefile{dataW}.\\
\textbf{Reference} & Documenta cada endpoint con request, parámetros, body, respuesta, errores e interpretación.\\
\bottomrule
\end{longtable}

\subsection{Convención de bloques de API}

Cada endpoint se documenta con la siguiente plantilla:

\begin{itemize}
    \item \textbf{Endpoint}: método HTTP y ruta.
    \item \textbf{Descripción}: propósito funcional.
    \item \textbf{Autenticación}: si requiere token y qué rol necesita.
    \item \textbf{Request}: path params, query params, headers y body.
    \item \textbf{Response}: códigos esperados y estructura de salida.
    \item \textbf{Errores}: significado de cada caso de fallo.
    \item \textbf{Interacción interna}: servicios que toca: Keycloak, DAO, Redis, Elasticsearch, PostgreSQL, MongoDB, Hive o Neo4J.
\end{itemize}

\section{Objetivos del sistema}

\subsection{Objetivo general}

Construir una plataforma distribuida de gestión de restaurantes que permita administrar entidades operacionales, autenticar usuarios por roles, ofrecer búsqueda textual eficiente, cachear datos consultados con frecuencia y generar análisis OLAP mediante un flujo de datos completo.

\subsection{Objetivos específicos}

\begin{enumerate}
    \item Exponer una API REST segura para autenticación, usuarios, restaurantes, menús, mesas, pedidos y reservaciones.
    \item Permitir el uso intercambiable de PostgreSQL o MongoDB como motor activo de persistencia.
    \item Implementar MongoDB en modo shardeado con config servers, shard replica set y router \codefile{mongos}.
    \item Integrar Keycloak para registro, login, emisión y validación de tokens JWT.
    \item Utilizar Redis para reducir accesos repetidos a la base de datos y a Elasticsearch.
    \item Utilizar Elasticsearch para búsqueda textual sobre productos del menú.
    \item Publicar servicios detrás de Nginx mediante rutas \codefile{/api} y \codefile{/search}.
    \item Construir un Data Warehouse OLAP con esquema estrella en Hive.
    \item Procesar datos con Spark usando DataFrames y SparkSQL.
    \item Visualizar indicadores en Superset.
    \item Orquestar el pipeline analítico con Airflow.
    \item Modelar y validar asignación de rutas con Neo4J y una heurística de vecino más cercano.
\end{enumerate}

\section{Fuentes documentales y artefactos base}

Esta documentación consolida información técnica del repositorio y de los artefactos ubicados en \codefile{documentation/}. Los documentos y diagramas base son:

\begin{itemize}
    \item \codefile{documentation/\_BD2\_PY01\_\_Restaurantes.pdf}: enunciado/guía de la primera etapa del proyecto.
    \item \codefile{documentation/\_BD2\_PY02\_\_Restaurantes\_2.pdf}: enunciado/guía de la segunda etapa del proyecto.
    \item \codefile{documentation/\_BD2\_TC01b\_\_Docker\_y\_RestAPI.pdf}: material de referencia sobre Docker y REST API.
    \item \codefile{documentation/Restaurantes e2 - Cesar Gamboa, Daniel Morera.pdf}: documento previo de entrega.
    \item \codefile{documentation/Diagrama de arquitectura física.png}: diagrama físico incluido en la sección de arquitectura física.
    \item \codefile{documentation/Diagrama de arquitectura lógica.png}: diagrama lógico incluido en la sección de arquitectura lógica.
    \item \codefile{documentation/Diagrama de flujo de datos - Admin\_Menu.png}: flujo de creación de menú por administrador.
    \item \codefile{documentation/Diagrama de flujo de datos - Cliente\_Búsqueda.png}: flujo de búsqueda de productos por cliente.
    \item \codefile{documentation/Diagrama de flujo de datos - Cliente\_Reserva.png}: flujo de reservación por cliente.
\end{itemize}

\section{Getting started}

Esta sección permite levantar el sistema y realizar una primera llamada autenticada. La idea es que cualquier persona nueva en el proyecto pueda pasar de cero a una petición funcional sin leer todavía toda la arquitectura.

\subsection{Requisitos previos}

\begin{itemize}
    \item Docker Desktop con Docker Compose.
    \item PowerShell en Windows.
    \item Puertos disponibles: \codefile{5432}, \codefile{6379}, \codefile{8000}, \codefile{8001}, \codefile{8080}, \codefile{9200}, \codefile{27017}, \codefile{8088}, \codefile{8090}, \codefile{10000}.
    \item Archivo \codefile{.env} en la raíz del proyecto.
\end{itemize}

\subsection{1. Levantar servicios}

\begin{lstlisting}[language=bash]
docker compose up -d --build
\end{lstlisting}

\subsection{2. Revisar estado}

\begin{lstlisting}[language=bash]
docker compose ps
docker compose logs -f keycloak
docker compose logs -f api
\end{lstlisting}

\subsection{3. Abrir documentación interactiva}

\begin{itemize}
    \item API directa: \codefile{http://localhost:8000/docs}
    \item API por Nginx: \codefile{http://localhost:8080/api/docs}
    \item Search por Nginx: \codefile{http://localhost:8080/search/docs}
\end{itemize}

\subsection{4. Obtener token}

\begin{lstlisting}[language=bash]
curl -X POST http://localhost:8000/auth/login ^
  -H "Content-Type: application/json" ^
  -d "{\"username\":\"admin\",\"password\":\"admin\"}"
\end{lstlisting}

La respuesta incluye \codefile{access\_token}. Ese token se usa como:

\begin{lstlisting}
Authorization: Bearer <access_token>
\end{lstlisting}

\subsection{5. Hacer primera llamada autenticada}

\begin{lstlisting}[language=bash]
curl http://localhost:8000/users/me ^
  -H "Authorization: Bearer <access_token>"
\end{lstlisting}

Si el token es válido y el usuario existe en la base activa, la API responde con los datos locales del usuario.

\section{Conceptos principales}

\subsection{Access token}

El \codefile{access\_token} es un JWT emitido por Keycloak. La API no acepta usuario y contraseña en cada endpoint; acepta un token firmado. El token identifica al usuario mediante el claim \codefile{sub} y contiene roles para autorización.

\subsection{Roles}

\begin{itemize}
    \item \codefile{client}: usuario cliente. Puede consultar datos protegidos, buscar productos, crear pedidos y crear reservaciones.
    \item \codefile{admin}: usuario administrador. Puede crear, actualizar y eliminar restaurantes, menús y mesas.
\end{itemize}

\subsection{Base URL}

La misma API puede consumirse por dos caminos:

\begin{itemize}
    \item \textbf{Directo}: \codefile{http://localhost:8000}
    \item \textbf{Proxy Nginx}: \codefile{http://localhost:8080/api}
\end{itemize}

El servicio de búsqueda se publica por Nginx como:

\begin{itemize}
    \item \codefile{http://localhost:8080/search}
\end{itemize}

\subsection{Motor activo de base de datos}

La API usa un motor activo a la vez. La variable \codefile{DATABASE\_ENGINE} define si las operaciones se ejecutan sobre PostgreSQL o MongoDB. Los routers no cambian; lo que cambia es la implementación del DAO.

\subsection{Cache hit y cache miss}

Cuando se consulta una lista de menús o una búsqueda, Redis puede devolver una respuesta guardada:

\begin{itemize}
    \item \textbf{Cache hit}: Redis tenía el resultado y se evita consultar base o Elasticsearch.
    \item \textbf{Cache miss}: Redis no tenía el resultado; la API consulta la fuente real y luego guarda la respuesta.
\end{itemize}

\subsection{Reindexado}

El endpoint \codefile{POST /reindex} reconstruye el índice de Elasticsearch desde los menús de la base activa. Debe ejecutarse cuando hay cambios de catálogo que no fueron indexados automáticamente o cuando Airflow detecta cambios en \codefile{dim\_product}.

\subsection{Data Warehouse}

El Data Warehouse no reemplaza la base transaccional. Es una copia analítica en Hive optimizada para agregaciones, cubos y dashboards.

\section{Estructura general del repositorio}

\begin{longtable}{p{0.25\textwidth}p{0.68\textwidth}}
\toprule
\textbf{Ruta} & \textbf{Propósito}\\
\midrule
\codefile{backend/} & Código Python de la API FastAPI, modelos SQLAlchemy, esquemas Pydantic, routers, DAO, autenticación, caché y búsqueda.\\
\codefile{dbs/postgres/} & Scripts de inicialización de PostgreSQL.\\
\codefile{dbs/mongo/} & Scripts de inicialización del cluster shardeado MongoDB, colecciones, índices y shard keys.\\
\codefile{nginx/} & Configuración de proxy reverso Nginx.\\
\codefile{dataW/} & Plataforma analítica: Data Warehouse, Spark, Superset, Airflow, pruebas y evidencias.\\
\codefile{neo4j/} & Carga de grafo, consultas Cypher y asignación de rutas de entrega.\\
\codefile{kubernetes/} & Manifiestos YAML y scripts PowerShell para despliegue en Kubernetes local.\\
\codefile{documentation/} & Diagramas, enunciados PDF y esta documentación LaTeX.\\
\codefile{docker-compose.yml} & Orquestación local de servicios transaccionales, analíticos y de grafo.\\
\codefile{README.md} & Guía rápida de ejecución, URLs, credenciales y pruebas.\\
\bottomrule
\end{longtable}

\section{Arquitectura lógica}

La arquitectura lógica representa las capas del sistema y sus responsabilidades.

\begin{figure}[H]
    \centering
    \includegraphics[width=\textwidth]{documentation/Diagrama de arquitectura lógica.png}
    \caption{Arquitectura lógica del sistema.}
\end{figure}

\subsection{Capas}

\begin{enumerate}
    \item \textbf{Clientes}: aplicaciones web o clientes HTTP que consumen endpoints REST.
    \item \textbf{Capa de presentación}: endpoints HTTP expuestos por FastAPI y publicados por Nginx.
    \item \textbf{Capa de aplicación}: routers, validaciones de negocio, autenticación, autorización, caché y búsqueda.
    \item \textbf{Capa de acceso a datos}: DAO común con implementaciones para PostgreSQL y MongoDB.
    \item \textbf{Modelos de datos}: entidades transaccionales como usuarios, roles, restaurantes, menús, mesas, pedidos, ítems y reservaciones.
    \item \textbf{Servicios externos}: Keycloak, Redis y Elasticsearch.
    \item \textbf{Capa analítica}: Hive, Spark, Superset y Airflow dentro de \codefile{dataW}.
\end{enumerate}

\section{Arquitectura física}

La arquitectura física define contenedores, red interna, puertos, servicios externos expuestos y volúmenes persistentes. Todos los contenedores principales se conectan a la red Docker \codefile{restaurant-network}.

\begin{figure}[H]
    \centering
    \includegraphics[width=\textwidth]{documentation/Diagrama de arquitectura física.png}
    \caption{Arquitectura física del sistema.}
\end{figure}

\subsection{Componentes físicos}

Los puertos de la siguiente tabla corresponden a los valores por defecto de \codefile{docker-compose.yml}. Si se sobrescriben variables en \codefile{.env}, los puertos locales pueden cambiar. Cuando se indica \textit{interno}, significa que el servicio escucha dentro de la red Docker y normalmente se consume a través de Nginx u otro contenedor, no directamente desde el host.

\begin{longtable}{p{0.22\textwidth}p{0.18\textwidth}p{0.48\textwidth}}
\toprule
\textbf{Servicio} & \textbf{Puerto local} & \textbf{Función}\\
\midrule
Nginx & \codefile{8080} & Entrada HTTP unificada. Redirige \codefile{/api/} hacia API y \codefile{/search/} hacia Search.\\
API FastAPI & \codefile{8000} directo, interno \codefile{80} & Servicio transaccional principal.\\
Search FastAPI & interno \codefile{80} & Microservicio especializado para búsqueda y reindexado.\\
Keycloak & \codefile{8001} & Identidad, roles, emisión de JWT y consola administrativa.\\
PostgreSQL & \codefile{5432} & Base relacional alternativa para datos transaccionales.\\
MongoDB \codefile{mongos} & \codefile{27017} & Router del cluster shardeado MongoDB.\\
Redis & \codefile{6379} & Caché de menús y búsquedas.\\
Elasticsearch & \codefile{9200} & Índice textual de productos.\\
Spark Master & \codefile{7077}, UI \codefile{8081} & Coordinador de procesamiento Spark.\\
Spark Worker & UI \codefile{8082} & Ejecutor de tareas Spark.\\
HiveServer2 & \codefile{10000}, UI \codefile{10002} & Data Warehouse OLAP.\\
Superset & \codefile{8088} & Visualización BI y dashboards.\\
Airflow & \codefile{8090} & Orquestación del pipeline analítico.\\
Neo4J & \codefile{7474}, \codefile{7687} & Grafo y rutas de entrega.\\
\bottomrule
\end{longtable}

\subsection{Volúmenes persistentes}

El sistema utiliza volúmenes Docker para conservar datos aunque se reinicien contenedores. Los más relevantes son:

\begin{itemize}
    \item \codefile{postgres\_data}: datos de PostgreSQL.
    \item \codefile{cfg1\_data}, \codefile{cfg2\_data}, \codefile{cfg3\_data}: config servers de MongoDB.
    \item \codefile{shard1a\_data}, \codefile{shard1b\_data}, \codefile{shard1c\_data}: shard replica set de MongoDB.
    \item \codefile{redis\_data}: persistencia de Redis.
    \item \codefile{elasticsearch\_data}: índice de Elasticsearch.
    \item \codefile{hive\_data} y \codefile{hive\_metastore}: datos y metastore de Hive.
    \item \codefile{superset\_home}: metadata y configuración de Superset.
    \item \codefile{airflow\_postgres\_data}: metadata de Airflow.
    \item \codefile{neo4j\_data}: base de datos de Neo4J.
\end{itemize}

\section{Tecnologías utilizadas}

\subsection{FastAPI}

FastAPI expone los endpoints REST. Usa Pydantic para validar entradas y salidas, dependencias para inyectar autenticación y DAO, y documentación automática OpenAPI en \codefile{/docs}. El proyecto usa dos aplicaciones principales:

\begin{itemize}
    \item \codefile{backend/app/api\_main.py}: API transaccional.
    \item \codefile{backend/app/search\_main.py}: microservicio de búsqueda.
\end{itemize}

\subsection{Keycloak}

Keycloak administra usuarios, roles y tokens JWT. El login se realiza contra el endpoint OpenID Connect de Keycloak usando \codefile{grant\_type=password}. Las rutas protegidas reciben el token en:

\begin{lstlisting}
Authorization: Bearer <access_token>
\end{lstlisting}

La API valida el JWT consultando las llaves públicas JWKS del realm configurado. Si el token es válido, el payload se usa para identificar al usuario y revisar roles. Los roles esperados son \codefile{admin} y \codefile{client}.

\subsection{PostgreSQL}

PostgreSQL funciona como motor relacional. Los modelos SQLAlchemy definen tablas, claves foráneas, relaciones y restricciones únicas. Se usa cuando:

\begin{lstlisting}
DATABASE_ENGINE=postgres
\end{lstlisting}

\subsection{MongoDB shardeado}

MongoDB funciona como motor documental distribuido. El cluster se compone de:

\begin{itemize}
    \item Tres config servers: \codefile{cfg1}, \codefile{cfg2}, \codefile{cfg3}.
    \item Un shard replica set: \codefile{shard1a}, \codefile{shard1b}, \codefile{shard1c}.
    \item Un router: \codefile{mongos}.
\end{itemize}

La API no se conecta a los shards directamente, sino a:

\begin{lstlisting}
mongodb://mongos:27017/restaurant_mongo_db
\end{lstlisting}

Las colecciones shardeadas principales son:

\begin{itemize}
    \item \codefile{menus}: shard key \codefile{\{ restaurant\_id: 1, dish\_name: 1 \}}.
    \item \codefile{reservations}: shard key \codefile{\{ reservation\_id: "hashed" \}}.
\end{itemize}

\subsection{Redis}

Redis guarda respuestas frecuentes para reducir latencia:

\begin{itemize}
    \item Lista completa de menús: \codefile{menus:all}.
    \item Menú individual: \codefile{menus:<id>}.
    \item Búsquedas textuales: \codefile{search:products:text:<query>}.
    \item Búsquedas por categoría: \codefile{search:products:category:<category>}.
\end{itemize}

El TTL general se controla con \codefile{CACHE\_TTL\_SECONDS}; el TTL de búsqueda se controla con \codefile{SEARCH\_CACHE\_TTL\_SECONDS}.

\subsection{Elasticsearch}

Elasticsearch indexa productos del menú en el índice \codefile{products}. Cada documento contiene:

\begin{itemize}
    \item \codefile{menu\_id}
    \item \codefile{dish\_name}
    \item \codefile{category}
    \item \codefile{description}
    \item \codefile{price}
    \item \codefile{restaurant\_id}
\end{itemize}

Las búsquedas usan \codefile{multi\_match} sobre nombre, categoría y descripción. El filtro por categoría usa un campo keyword normalizado en minúsculas: \codefile{category.raw}.

\subsection{Nginx}

Nginx opera como proxy reverso. Publica:

\begin{itemize}
    \item \codefile{http://localhost:8080/api/} hacia el servicio \codefile{api}.
    \item \codefile{http://localhost:8080/search/} hacia el servicio \codefile{search}.
\end{itemize}

\subsection{Docker Compose}

Docker Compose levanta el ambiente completo. Además de los servicios transaccionales, integra Spark, Hive, Superset, Airflow y Neo4J. El comando base es:

\begin{lstlisting}[language=bash]
docker compose up -d --build
\end{lstlisting}

\subsection{Kubernetes}

La carpeta \codefile{kubernetes/} contiene manifiestos para namespace, configmaps, servicios, deployments, storage e ingress. Su objetivo es reproducir los servicios principales fuera de Compose, usando scripts PowerShell de inicialización.

\subsection{Apache Hive}

Hive implementa el Data Warehouse. El esquema estrella vive en \codefile{dataW/warehouse/schemas/schema\_star.sql}. HiveServer2 permite que Superset y validaciones consulten la base \codefile{restaurant\_warehouse}.

\subsection{Apache Spark}

Spark procesa datos a escala. El job \codefile{dataW/spark/jobs/analytics\_job.py} crea DataFrames, aplica transformaciones, ejecuta SparkSQL y genera salidas CSV/JSON.

\subsection{Apache Superset}

Superset se conecta directamente a HiveServer2 mediante la conexión:

\begin{lstlisting}
hive://hive@hiveserver2:10000/restaurant_warehouse?auth=NOSASL
\end{lstlisting}

Su función es construir dashboards sobre vistas OLAP.

\subsection{Apache Airflow}

Airflow orquesta el pipeline analítico diario \codefile{restaurant\_olap\_pipeline}. Coordina extracción, Spark, carga a Hive, validación OLAP y reindexado opcional de Elasticsearch.

\subsection{Neo4J}

Neo4J modela usuarios, pedidos, productos, restaurantes, ubicaciones y repartidores como grafo. El punto de rutas usa una heurística de vecino más cercano para asignar pedidos a repartidores disponibles.

\section{Configuración principal}

Las variables se leen desde \codefile{.env} y desde el bloque común de \codefile{docker-compose.yml}.

\begin{longtable}{p{0.33\textwidth}p{0.57\textwidth}}
\toprule
\textbf{Variable} & \textbf{Uso}\\
\midrule
\codefile{DATABASE\_ENGINE} & Selecciona \codefile{mongo} o \codefile{postgres}.\\
\codefile{POSTGRES\_URL} & URL SQLAlchemy de PostgreSQL.\\
\codefile{MONGO\_URL} & URL de conexión al router \codefile{mongos}.\\
\codefile{KEYCLOAK\_URL} & URL interna o externa de Keycloak.\\
\codefile{KEYCLOAK\_REALM} & Realm de Keycloak: \codefile{restaurant-realm}.\\
\codefile{KEYCLOAK\_CLIENT\_ID} & Cliente OIDC usado por la API.\\
\codefile{REDIS\_URL} & URL de Redis.\\
\codefile{ELASTICSEARCH\_URL} & URL de Elasticsearch.\\
\codefile{CACHE\_TTL\_SECONDS} & TTL de caché general.\\
\codefile{SEARCH\_CACHE\_TTL\_SECONDS} & TTL de caché de búsqueda.\\
\codefile{SPARK\_SCALE\_FACTOR} & Multiplicador de volumen para el job Spark.\\
\codefile{AIRFLOW\_PORT} & Puerto de Airflow Webserver.\\
\codefile{SUPERSET\_PORT} & Puerto de Superset.\\
\codefile{ENABLE\_DELIVERY\_ROUTE\_VALIDATION} & Activa validación opcional de rutas en Airflow.\\
\bottomrule
\end{longtable}

\section{Modelo de datos transaccional}

\subsection{Entidades}

\begin{longtable}{p{0.20\textwidth}p{0.70\textwidth}}
\toprule
\textbf{Entidad} & \textbf{Campos principales}\\
\midrule
\codefile{Role} & \codefile{role\_id}, \codefile{role\_name}.\\
\codefile{User} & \codefile{user\_id}, \codefile{user\_name}, \codefile{keycloak\_id}, \codefile{role\_id}.\\
\codefile{Restaurant} & \codefile{restaurant\_id}, \codefile{restaurant\_name}, \codefile{admin\_id}, \codefile{restaurant\_status}.\\
\codefile{Menu} & \codefile{menu\_id}, \codefile{dish\_name}, \codefile{category}, \codefile{description}, \codefile{price}, \codefile{restaurant\_id}.\\
\codefile{Table} & \codefile{table\_id}, \codefile{table\_number}, \codefile{table\_status}, \codefile{restaurant\_id}.\\
\codefile{Order} & \codefile{order\_id}, \codefile{table\_id}, \codefile{client\_id}, \codefile{order\_type}, \codefile{restaurant\_id}.\\
\codefile{OrderItem} & \codefile{order\_item\_id}, \codefile{order\_id}, \codefile{menu\_id}, \codefile{quantity}, \codefile{price}.\\
\codefile{Reservation} & \codefile{reservation\_id}, \codefile{table\_id}, \codefile{client\_id}, \codefile{reservation\_date}, \codefile{reservation\_status}.\\
\bottomrule
\end{longtable}

\subsection{Restricciones relevantes}

\begin{itemize}
    \item En menús no puede repetirse el mismo \codefile{dish\_name} dentro del mismo restaurante.
    \item En mesas no puede repetirse el mismo \codefile{table\_number} dentro del mismo restaurante.
    \item \codefile{keycloak\_id} identifica al usuario en Keycloak y debe ser único si existe.
    \item \codefile{table\_status} usa valores entre 0 y 2: ocupado, disponible o reservado, según la convención del código.
    \item Las operaciones de creación validan claves foráneas lógicas antes de persistir.
\end{itemize}

\section{Capa DAO}

La API usa \codefile{get\_dao()} para seleccionar una implementación concreta:

\begin{itemize}
    \item \codefile{PostgresDAO}: usa SQLAlchemy y sesiones transaccionales.
    \item \codefile{MongoDAO}: usa PyMongo, colecciones y contadores para generar IDs incrementales.
\end{itemize}

Esto permite que los routers no conozcan el motor activo. Por ejemplo, \codefile{dao.create\_menu(...)} crea un documento en MongoDB o una fila en PostgreSQL según \codefile{DATABASE\_ENGINE}.

\begin{lstlisting}[language=Python,caption={Selección conceptual del DAO activo.}]
if DATABASE_ENGINE == "mongo":
    yield MongoDAO(get_mongo_db())
elif DATABASE_ENGINE == "postgres":
    yield PostgresDAO(db)
\end{lstlisting}

\section{Autenticación y autorización}

\subsection{Flujo de registro}

El registro crea primero el usuario en Keycloak con rol \codefile{client}. Luego crea el usuario local en la base activa. Si la creación local falla, se intenta revertir eliminando el usuario creado en Keycloak.

\subsection{Flujo de login}

El login envía usuario y contraseña al endpoint OIDC de Keycloak. Si las credenciales son válidas, Keycloak responde con un JSON que incluye \codefile{access\_token}, \codefile{refresh\_token}, expiración y metadatos del token.

\subsection{Validación de token}

Cada ruta protegida usa \codefile{HTTPBearer}. La API:

\begin{enumerate}
    \item Extrae el token del header \codefile{Authorization}.
    \item Consulta las llaves públicas JWKS de Keycloak.
    \item Busca la llave con el \codefile{kid} del token.
    \item Decodifica el JWT con algoritmo \codefile{RS256}.
    \item Valida el \codefile{issuer}.
    \item Devuelve el payload si todo es correcto.
\end{enumerate}

\subsection{Roles}

\begin{itemize}
    \item \textbf{client}: puede consultar recursos protegidos, crear pedidos y reservas.
    \item \textbf{admin}: puede crear, actualizar y eliminar restaurantes, menús y mesas; también puede administrar usuarios cuando la lógica del endpoint lo permite.
\end{itemize}

\section{Convenciones generales de API}

\subsection{Base URLs}

\begin{itemize}
    \item API por Nginx: \codefile{http://localhost:8080/api}
    \item Search por Nginx: \codefile{http://localhost:8080/search}
    \item API directa: \codefile{http://localhost:8000}
    \item Documentación API directa: \codefile{http://localhost:8000/docs}
    \item Documentación por Nginx: \codefile{http://localhost:8080/api/docs}
\end{itemize}

\subsection{Headers}

Todas las rutas protegidas requieren:

\begin{lstlisting}
Authorization: Bearer <access_token>
Content-Type: application/json
\end{lstlisting}

\subsection{Códigos de estado}

\begin{longtable}{p{0.18\textwidth}p{0.72\textwidth}}
\toprule
\textbf{Código} & \textbf{Interpretación}\\
\midrule
\codefile{200} & Operación exitosa con cuerpo de respuesta.\\
\codefile{201} & Recurso creado correctamente.\\
\codefile{204} & Operación exitosa sin cuerpo, normalmente eliminación.\\
\codefile{400} & Solicitud inválida o regla de negocio incumplida.\\
\codefile{401} & Token ausente, inválido, expirado o credenciales incorrectas.\\
\codefile{403} & Usuario autenticado sin permisos suficientes.\\
\codefile{404} & Recurso o dependencia no encontrada.\\
\codefile{409} & Conflicto de unicidad, por ejemplo plato o mesa duplicada.\\
\codefile{422} & Error de validación Pydantic: tipos, rangos o campos requeridos.\\
\codefile{500} & Error interno no controlado.\\
\codefile{503} & Servicio externo no disponible o mal configurado, especialmente Keycloak.\\
\bottomrule
\end{longtable}

\section{How-Tos}

Los \textit{How-Tos} son guías cortas orientadas a tareas. A diferencia de la referencia, aquí importa el flujo completo.

\subsection{Registrar un cliente}

\begin{enumerate}
    \item Enviar \endpoint{POST}{/auth/register} con \codefile{username}, \codefile{email} y \codefile{password}.
    \item La API crea el usuario en Keycloak con rol \codefile{client}.
    \item La API crea el usuario local en PostgreSQL o MongoDB.
    \item La respuesta devuelve \codefile{user\_id} y \codefile{username}.
\end{enumerate}

\subsection{Crear un plato de menú}

\begin{enumerate}
    \item Iniciar sesión como administrador usando \endpoint{POST}{/auth/login}.
    \item Enviar \endpoint{POST}{/menus/} con el token admin.
    \item La API valida rol, restaurante existente y unicidad de plato por restaurante.
    \item La API guarda el menú, borra \codefile{menus:all} y borra cachés \codefile{search:products:*}.
\end{enumerate}

\subsection{Buscar productos}

\begin{enumerate}
    \item Iniciar sesión como cliente o admin.
    \item Ejecutar \endpoint{GET}{/products?q=pizza} en el servicio Search.
    \item Redis se consulta primero.
    \item Si no existe caché, Elasticsearch ejecuta \codefile{multi\_match}.
    \item La respuesta se guarda en Redis y se devuelve al cliente.
\end{enumerate}

\subsection{Reindexar Elasticsearch}

\begin{enumerate}
    \item Enviar \endpoint{POST}{/reindex} con token válido.
    \item El servicio Search recrea el índice \codefile{products}.
    \item Consulta todos los menús desde el DAO activo.
    \item Indexa cada producto y refresca el índice.
    \item Elimina caché de búsqueda.
\end{enumerate}

\subsection{Ejecutar el pipeline analítico}

\begin{enumerate}
    \item Levantar Hive, Spark, Airflow y Superset con Docker Compose.
    \item Abrir \codefile{http://localhost:8090}.
    \item Activar y disparar el DAG \codefile{restaurant\_olap\_pipeline}.
    \item Revisar \codefile{Grid} y logs de cada tarea.
    \item Abrir Superset en \codefile{http://localhost:8088} para consultar dashboards.
\end{enumerate}

\section{API Reference}

Esta sección usa el estilo de referencia de una documentación de plataforma: cada operación indica método, ruta, autenticación, parámetros, body, respuesta e interacción interna.

\subsection{Auth}

\subsubsection{\methodpost{} \codefile{/auth/register}}

\textbf{Descripción}: registra un cliente nuevo.

\textbf{Autenticación}: no requiere token.

\textbf{Request body}
\begin{longtable}{p{0.20\textwidth}p{0.18\textwidth}p{0.15\textwidth}p{0.37\textwidth}}
\toprule
\textbf{Campo} & \textbf{Tipo} & \textbf{Requerido} & \textbf{Descripción}\\
\midrule
\codefile{username} & string & Sí & 3 a 64 caracteres. Nombre único local.\\
\codefile{email} & string & Sí & Email válido para Keycloak.\\
\codefile{password} & string & Sí & 8 a 128 caracteres.\\
\bottomrule
\end{longtable}

\textbf{Response 201}
\begin{lstlisting}[language=json]
{
  "message": "Usuario creado correctamente",
  "user_id": 2,
  "username": "cliente1"
}
\end{lstlisting}

\textbf{Errores}: \codefile{400} usuario duplicado; \codefile{503} Keycloak no configurado; \codefile{500} fallo creando usuario.

\textbf{Interacción interna}: Keycloak Admin API, DAO activo, tabla/colección \codefile{users}, tabla/colección \codefile{roles}.

\subsubsection{\methodpost{} \codefile{/auth/login}}

\textbf{Descripción}: autentica credenciales contra Keycloak y devuelve tokens.

\textbf{Autenticación}: no requiere token.

\textbf{Request body}
\begin{longtable}{p{0.20\textwidth}p{0.18\textwidth}p{0.15\textwidth}p{0.37\textwidth}}
\toprule
\textbf{Campo} & \textbf{Tipo} & \textbf{Requerido} & \textbf{Descripción}\\
\midrule
\codefile{username} & string & Sí & Usuario de Keycloak.\\
\codefile{password} & string & Sí & Contraseña del usuario.\\
\bottomrule
\end{longtable}

\textbf{Response 200}: objeto OIDC de Keycloak con \codefile{access\_token}, \codefile{refresh\_token}, expiración, tipo de token y scope.

\textbf{Errores}: \codefile{401} credenciales inválidas; \codefile{403} acciones pendientes en Keycloak.

\textbf{Interacción interna}: endpoint OIDC \codefile{/realms/<realm>/protocol/openid-connect/token}.

\subsection{Users}

\subsubsection{\methodget{} \codefile{/users/me}}

\textbf{Descripción}: devuelve el usuario local asociado al token actual.

\textbf{Autenticación}: token válido.

\textbf{Headers}: \codefile{Authorization: Bearer <access\_token>}.

\textbf{Response 200}
\begin{lstlisting}[language=json]
{
  "user_name": "admin",
  "role_id": 1,
  "user_id": 1
}
\end{lstlisting}

\textbf{Errores}: \codefile{401} token sin \codefile{sub}; \codefile{404} usuario no existe en base local.

\textbf{Interacción interna}: validación JWKS contra Keycloak y consulta \codefile{get\_user\_by\_keycloak\_id}.

\subsubsection{\methodput{} \codefile{/users/{user_id}}}

\textbf{Descripción}: actualiza el perfil de un usuario.

\textbf{Autenticación}: dueño del recurso o admin.

\textbf{Path params}
\begin{longtable}{p{0.20\textwidth}p{0.18\textwidth}p{0.52\textwidth}}
\toprule
\textbf{Parámetro} & \textbf{Tipo} & \textbf{Descripción}\\
\midrule
\codefile{user\_id} & integer & Identificador del usuario local.\\
\bottomrule
\end{longtable}

\textbf{Request body}
\begin{lstlisting}[language=json]
{
  "user_name": "nuevo_nombre"
}
\end{lstlisting}

\textbf{Response 200}: \codefile{UserResponse}.

\textbf{Errores}: \codefile{400} username ocupado; \codefile{403} no es dueño/admin o intenta cambiar rol; \codefile{404} usuario inexistente; \codefile{503} no se pudo sincronizar Keycloak.

\textbf{Interacción interna}: Keycloak Admin API y DAO activo.

\subsubsection{\methoddelete{} \codefile{/users/{user_id}}}

\textbf{Descripción}: elimina cuenta local y usuario de Keycloak.

\textbf{Autenticación}: dueño del recurso o admin.

\textbf{Response 204}: sin body.

\textbf{Errores}: \codefile{401}, \codefile{403}, \codefile{404}, \codefile{503}.

\textbf{Interacción interna}: Keycloak Admin API y DAO activo.

\subsection{Restaurants}

\subsubsection{\methodget{} \codefile{/restaurants/}}

\textbf{Descripción}: lista restaurantes.

\textbf{Autenticación}: token válido.

\textbf{Response 200}: arreglo de \codefile{RestaurantResponse}.

\textbf{Interacción interna}: DAO activo.

\subsubsection{\methodget{} \codefile{/restaurants/{restaurant_id}}}

\textbf{Descripción}: obtiene un restaurante por ID.

\textbf{Autenticación}: token válido.

\textbf{Response 200}
\begin{lstlisting}[language=json]
{
  "restaurant_name": "Restaurante Central",
  "admin_id": 1,
  "restaurant_status": 1,
  "restaurant_id": 1
}
\end{lstlisting}

\textbf{Errores}: \codefile{404} restaurante no encontrado.

\subsubsection{\methodpost{} \codefile{/restaurants/}}

\textbf{Descripción}: crea restaurante.

\textbf{Autenticación}: admin.

\textbf{Request body}
\begin{lstlisting}[language=json]
{
  "restaurant_name": "Restaurante Central",
  "admin_id": 1,
  "restaurant_status": 1
}
\end{lstlisting}

\textbf{Response 201}: \codefile{RestaurantResponse}.

\textbf{Errores}: \codefile{403} no admin; \codefile{404} administrador no existe; \codefile{422} body inválido.

\textbf{Interacción interna}: validación de rol, DAO activo, relación con \codefile{users}.

\subsubsection{\methodput{} \codefile{/restaurants/{restaurant_id}}}

\textbf{Descripción}: actualiza parcialmente restaurante.

\textbf{Autenticación}: admin.

\textbf{Request body}: cualquiera de \codefile{restaurant\_name}, \codefile{admin\_id}, \codefile{restaurant\_status}.

\textbf{Errores}: \codefile{400} body vacío; \codefile{403} no admin; \codefile{404} restaurante o admin no existe.

\subsubsection{\methoddelete{} \codefile{/restaurants/{restaurant_id}}}

\textbf{Descripción}: elimina restaurante.

\textbf{Autenticación}: admin.

\textbf{Response 204}: sin body.

\textbf{Errores}: \codefile{403} no admin; \codefile{404} restaurante no encontrado.

\subsection{Menus}

\subsubsection{\methodget{} \codefile{/menus/}}

\textbf{Descripción}: lista platos del menú.

\textbf{Autenticación}: token válido.

\textbf{Response 200}: arreglo de \codefile{MenuResponse}.

\textbf{Cache}: lee/escribe \codefile{menus:all}.

\textbf{Interacción interna}: Redis y DAO activo.

\subsubsection{\methodget{} \codefile{/menus/\{menu\_id\}}}

\textbf{Descripción}: obtiene un plato por ID.

\textbf{Autenticación}: token válido.

\textbf{Response 200}
\begin{lstlisting}[language=json]
{
  "dish_name": "Pizza Margarita",
  "category": "Plato fuerte",
  "description": "Pizza con queso, tomate y albahaca",
  "price": 7200.0,
  "restaurant_id": 1,
  "menu_id": 1
}
\end{lstlisting}

\textbf{Cache}: lee/escribe \codefile{menus:<menu\_id>}.

\textbf{Errores}: \codefile{404} plato no encontrado.

\subsubsection{\methodpost{} \codefile{/menus/}}

\textbf{Descripción}: crea plato.

\textbf{Autenticación}: admin.

\textbf{Request body}
\begin{lstlisting}[language=json]
{
  "dish_name": "Pizza Margarita",
  "category": "Plato fuerte",
  "description": "Pizza con queso, tomate y albahaca",
  "price": 7200.0,
  "restaurant_id": 1
}
\end{lstlisting}

\textbf{Response 201}: \codefile{MenuResponse}.

\textbf{Errores}: \codefile{403} no admin; \codefile{404} restaurante no existe; \codefile{409} plato duplicado en restaurante.

\textbf{Interacción interna}: DAO activo, Redis, invalidación de \codefile{menus:all} y \codefile{search:products:*}.

\subsubsection{\methodput{} \codefile{/menus/\{menu\_id\}}}

\textbf{Descripción}: actualiza parcialmente plato.

\textbf{Autenticación}: admin.

\textbf{Request body}: cualquiera de \codefile{dish\_name}, \codefile{category}, \codefile{description}, \codefile{price}, \codefile{restaurant\_id}.

\textbf{Errores}: \codefile{400} body vacío; \codefile{403}; \codefile{404}; \codefile{409}.

\textbf{Interacción interna}: DAO activo e invalidación de caché de menú y búsqueda.

\subsubsection{\methoddelete{} \codefile{/menus/\{menu\_id\}}}

\textbf{Descripción}: elimina plato.

\textbf{Autenticación}: admin.

\textbf{Response 204}: sin body.

\textbf{Interacción interna}: DAO activo e invalidación de \codefile{menus:all}, \codefile{menus:<id>} y \codefile{search:products:*}.

\subsection{Tables}

\subsubsection{\methodget{} \codefile{/tables/}}

\textbf{Descripción}: lista mesas.

\textbf{Autenticación}: token válido.

\textbf{Response 200}: arreglo de \codefile{TableResponse}.

\subsubsection{\methodpost{} \codefile{/tables/}}

\textbf{Descripción}: crea mesa.

\textbf{Autenticación}: admin.

\textbf{Request body}
\begin{lstlisting}[language=json]
{
  "table_number": 12,
  "table_status": 1,
  "restaurant_id": 1
}
\end{lstlisting}

\textbf{Response 201}: \codefile{TableResponse}.

\textbf{Errores}: \codefile{403} no admin; \codefile{404} restaurante no existe; \codefile{409} número de mesa duplicado en restaurante.

\subsubsection{\methodget{} \texttt{/tables/\{table\_id\}}}

\textbf{Descripción}: obtiene mesa por ID.

\textbf{Errores}: \codefile{404} mesa no encontrada.

\subsubsection{\methodput{} \texttt{/tables/\{table\_id\}}}

\textbf{Descripción}: actualiza mesa parcialmente.

\textbf{Autenticación}: admin.

\textbf{Errores}: \codefile{400} body vacío; \codefile{403}; \codefile{404}; \codefile{409}.

\subsubsection{\methoddelete{} \texttt{/tables/\{table\_id\}}}

\textbf{Descripción}: elimina mesa.

\textbf{Autenticación}: admin.

\textbf{Response 204}: sin body.

\subsection{Orders}

\subsubsection{\methodpost{} \codefile{/orders/}}

\textbf{Descripción}: crea pedido.

\textbf{Autenticación}: token válido.

\textbf{Request body}
\begin{lstlisting}[language=json]
{
  "table_id": 1,
  "client_id": 2,
  "order_type": "local",
  "restaurant_id": 1
}
\end{lstlisting}

\textbf{Response 201}: \codefile{OrderResponse}.

\textbf{Errores}: \codefile{404} mesa, cliente o restaurante no existe; \codefile{422} body inválido.

\textbf{Interacción interna}: DAO activo.

\subsubsection{\methodget{} \codefile{/orders/\{order\_id\}}}

\textbf{Descripción}: obtiene pedido por ID.

\textbf{Autenticación}: token válido.

\textbf{Errores}: \codefile{404} pedido no encontrado.

\subsection{Reservations}

\subsubsection{\methodpost{} \codefile{/reservations/}}

\textbf{Descripción}: crea reservación.

\textbf{Autenticación}: token válido.

\textbf{Request body}
\begin{lstlisting}[language=json]
{
  "table_id": 1,
  "client_id": 2,
  "reservation_date": "2026-06-25T19:30:00",
  "reservation_status": 1
}
\end{lstlisting}

\textbf{Response 201}: \codefile{ReservationResponse}.

\textbf{Errores}: \codefile{404} mesa o cliente no existe; \codefile{422} fecha o tipos inválidos.

\subsubsection{\methoddelete{} \codefile{/reservations/{reservation_id}}}

\textbf{Descripción}: cancela/elimina reservación.

\textbf{Autenticación}: token válido.

\textbf{Response 204}: sin body.

\textbf{Errores}: \codefile{404} reservación no encontrada.

\subsection{Search}

\subsubsection{\methodget{} \codefile{/products}}

\textbf{Descripción}: búsqueda textual de productos.

\textbf{Autenticación}: token válido.

\textbf{Query params}
\begin{longtable}{p{0.18\textwidth}p{0.16\textwidth}p{0.16\textwidth}p{0.40\textwidth}}
\toprule
\textbf{Parámetro} & \textbf{Tipo} & \textbf{Requerido} & \textbf{Descripción}\\
\midrule
\codefile{q} & string & Sí & Texto a buscar en nombre, categoría o descripción.\\
\bottomrule
\end{longtable}

\textbf{Response 200}: arreglo de productos indexados.

\textbf{Cache}: \codefile{search:products:text:<q>}.

\textbf{Interacción interna}: Redis y Elasticsearch.

\subsubsection{\methodget{} \codefile{/products/category/{category}}}

\textbf{Descripción}: busca productos por categoría exacta normalizada.

\textbf{Autenticación}: token válido.

\textbf{Response 200}: arreglo de productos indexados.

\textbf{Cache}: \codefile{search:products:category:<category>}.

\subsubsection{\methodpost{} \codefile{/reindex}}

\textbf{Descripción}: reconstruye índice de productos.

\textbf{Autenticación}: token válido.

\textbf{Response 200}
\begin{lstlisting}[language=json]
{
  "message": "Productos reindexados correctamente",
  "total": 25
}
\end{lstlisting}

\textbf{Interacción interna}: DAO activo, Elasticsearch y Redis.

\subsection{Graph}

\subsubsection{\methodget{} \codefile{/graph/export}}

\textbf{Descripción}: exporta datos operacionales para Neo4J y Data Warehouse.

\textbf{Autenticación}: en la implementación actual no declara dependencia de token.

\textbf{Response 200}
\begin{lstlisting}[language=json]
{
  "roles": [],
  "users": [],
  "restaurants": [],
  "products": [],
  "orders": [],
  "order_items": []
}
\end{lstlisting}

\textbf{Interacción interna}: DAO activo. También calcula \codefile{total} de cada pedido usando \codefile{order\_items}.

\section{Contratos de API}

\subsection{Sistema}

\begin{longtable}{p{0.25\textwidth}p{0.18\textwidth}p{0.45\textwidth}}
\toprule
\textbf{Endpoint} & \textbf{Auth} & \textbf{Respuesta}\\
\midrule
\endpoint{GET}{/ping} & No & \codefile{\{"message":"pong"\}}.\\
\endpoint{GET}{/} & No & Mensaje de API funcionando y versión.\\
\endpoint{GET}{/health} & No & \codefile{\{"status":"ok"\}}.\\
\bottomrule
\end{longtable}

\subsection{Autenticación}

\subsubsection{\endpoint{POST}{/auth/register}}

Registra un usuario cliente en Keycloak y en la base activa.

\textbf{Body}
\begin{lstlisting}[language=json]
{
  "username": "cliente1",
  "email": "cliente1@example.com",
  "password": "password123"
}
\end{lstlisting}

\textbf{Validaciones}
\begin{itemize}
    \item \codefile{username}: 3 a 64 caracteres.
    \item \codefile{email}: email válido.
    \item \codefile{password}: 8 a 128 caracteres.
    \item No debe existir otro usuario con el mismo \codefile{username}.
    \item Keycloak debe estar configurado y accesible.
\end{itemize}

\textbf{Respuesta 201}
\begin{lstlisting}[language=json]
{
  "message": "Usuario creado correctamente",
  "user_id": 2,
  "username": "cliente1"
}
\end{lstlisting}

\textbf{Errores}
\begin{itemize}
    \item \codefile{400}: campos faltantes o usuario duplicado.
    \item \codefile{503}: Keycloak no configurado.
    \item \codefile{500}: error creando usuario en Keycloak o en base de datos.
\end{itemize}

\subsubsection{\endpoint{POST}{/auth/login}}

Obtiene tokens desde Keycloak.

\textbf{Body}
\begin{lstlisting}[language=json]
{
  "username": "admin",
  "password": "admin"
}
\end{lstlisting}

\textbf{Respuesta 200}
\begin{lstlisting}[language=json]
{
  "access_token": "<jwt>",
  "expires_in": 300,
  "refresh_expires_in": 1800,
  "refresh_token": "<jwt>",
  "token_type": "Bearer",
  "not-before-policy": 0,
  "session_state": "...",
  "scope": "profile email"
}
\end{lstlisting}

\textbf{Interpretación}: \codefile{access\_token} se usa en rutas protegidas. \codefile{refresh\_token} permite renovar sesión según la configuración de Keycloak.

\textbf{Errores}
\begin{itemize}
    \item \codefile{401}: credenciales inválidas o usuario inexistente.
    \item \codefile{403}: acciones requeridas pendientes en Keycloak.
\end{itemize}

\subsection{Usuarios}

\begin{longtable}{p{0.20\textwidth}p{0.18\textwidth}p{0.20\textwidth}p{0.34\textwidth}}
\toprule
\textbf{Endpoint} & \textbf{Entrada} & \textbf{Auth} & \textbf{Salida / reglas}\\
\midrule
\endpoint{GET}{/users/me} & Ninguna & Token válido & Devuelve el usuario local asociado al \codefile{sub} del JWT.\\
\endpoint{PUT}{/users/{id}} & \codefile{UserUpdate} & Dueño o admin & Actualiza \codefile{user\_name}. No permite escalar rol. Sincroniza nombre con Keycloak.\\
\endpoint{DELETE}{/users/{id}} & ID en path & Dueño o admin & Elimina usuario en Keycloak y base local. Responde \codefile{204}.\\
\bottomrule
\end{longtable}

\textbf{UserUpdate}
\begin{lstlisting}[language=json]
{
  "user_name": "nuevo_nombre",
  "role_id": 2
}
\end{lstlisting}

Aunque el esquema acepta \codefile{role\_id}, el endpoint rechaza cambios de rol si difieren del rol actual.

\subsection{Restaurantes}

\small
\renewcommand{\arraystretch}{1.25}

\begin{tabularx}{\textwidth}{
>{\raggedright\arraybackslash}p{0.22\textwidth}
>{\raggedright\arraybackslash}p{0.22\textwidth}
>{\raggedright\arraybackslash}p{0.15\textwidth}
>{\raggedright\arraybackslash}X
}
\toprule
\textbf{Endpoint} & \textbf{Entrada} & \textbf{Auth} & \textbf{Comportamiento}\\
\midrule

\endpoint{GET}{/restaurants/}
& Ninguna
& Token válido
& Lista todos los restaurantes.\\

\endpoint{GET}{/restaurants/{id}}
& ID
& Token válido
& Devuelve un restaurante o responde 404.\\

\endpoint{POST}{/restaurants/}
& \codefile{RestaurantCreate}
& Admin
& Crea restaurante si \codefile{admin_id} existe.\\

\endpoint{PUT}{/restaurants/{id}}
& \codefile{RestaurantUpdate}
& Admin
& Actualiza parcialmente; requiere al menos un campo.\\

\endpoint{DELETE}{/restaurants/{id}}
& ID
& Admin
& Elimina restaurante o responde 404.\\

\bottomrule
\end{tabularx}

\normalsize
\renewcommand{\arraystretch}{1}

\vspace{0.6cm}

\textbf{RestaurantCreate}
\begin{lstlisting}[language=json]
{
  "restaurant_name": "Restaurante Central",
  "admin_id": 1,
  "restaurant_status": 1
}
\end{lstlisting}

\textbf{RestaurantResponse}
\begin{lstlisting}[language=json]
{
  "restaurant_name": "Restaurante Central",
  "admin_id": 1,
  "restaurant_status": 1,
  "restaurant_id": 1
}
\end{lstlisting}

\subsection{Menús}

\begin{longtable}{p{0.20\textwidth}p{0.22\textwidth}p{0.18\textwidth}p{0.32\textwidth}}
\toprule
\textbf{Endpoint} & \textbf{Entrada} & \textbf{Auth} & \textbf{Comportamiento}\\
\midrule
\endpoint{GET}{/menus/} & Ninguna & Token válido & Lista menús. Usa Redis con llave \codefile{menus:all}.\\
\endpoint{GET}{/menus/{id}} & ID & Token válido & Devuelve un plato. Usa Redis con llave \codefile{menus:<id>}.\\
\endpoint{POST}{/menus/} & \codefile{MenuCreate} & Admin & Crea plato si restaurante existe. Invalida caché.\\
\endpoint{PUT}{/menus/{id}} & \codefile{MenuUpdate} & Admin & Actualiza parcialmente. Invalida caché de menú y búsqueda.\\
\endpoint{DELETE}{/menus/{id}} & ID & Admin & Elimina plato e invalida caché.\\
\bottomrule
\end{longtable}

\textbf{MenuCreate}
\begin{lstlisting}[language=json]
{
  "dish_name": "Pizza Margarita",
  "category": "Plato fuerte",
  "description": "Pizza con queso, tomate y albahaca",
  "price": 7200.0,
  "restaurant_id": 1
}
\end{lstlisting}

\textbf{Reglas}
\begin{itemize}
    \item \codefile{price} debe ser mayor que 0.
    \item \codefile{restaurant\_id} debe existir.
    \item Si ya existe el mismo \codefile{dish\_name} en el restaurante, responde \codefile{409}.
\end{itemize}

\subsection{Mesas}

\begin{longtable}{p{0.20\textwidth}p{0.22\textwidth}p{0.18\textwidth}p{0.32\textwidth}}
\toprule
\textbf{Endpoint} & \textbf{Entrada} & \textbf{Auth} & \textbf{Comportamiento}\\
\midrule
\endpoint{GET}{/tables/} & Ninguna & Token válido & Lista mesas.\\
\endpoint{GET}{/tables/{id}} & ID & Token válido & Devuelve mesa o \codefile{404}.\\
\endpoint{POST}{/tables/} & \codefile{TableCreate} & Admin & Crea mesa si restaurante existe.\\
\endpoint{PUT}{/tables/{id}} & \codefile{TableUpdate} & Admin & Actualiza parcialmente.\\
\endpoint{DELETE}{/tables/{id}} & ID & Admin & Elimina mesa.\\
\bottomrule
\end{longtable}

\textbf{TableCreate}
\begin{lstlisting}[language=json]
{
  "table_number": 12,
  "table_status": 1,
  "restaurant_id": 1
}
\end{lstlisting}

\textbf{Reglas}
\begin{itemize}
    \item \codefile{table\_number} debe ser mayor que 0.
    \item \codefile{table\_status} debe estar entre 0 y 2.
    \item No se permite repetir número de mesa dentro del mismo restaurante.
\end{itemize}

\subsection{Pedidos}

\begin{longtable}{p{0.20\textwidth}p{0.22\textwidth}p{0.18\textwidth}p{0.32\textwidth}}
\toprule
\textbf{Endpoint} & \textbf{Entrada} & \textbf{Auth} & \textbf{Comportamiento}\\
\midrule
\endpoint{POST}{/orders/} & \codefile{OrderCreate} & Token válido & Crea pedido validando mesa opcional, cliente y restaurante.\\
\endpoint{GET}{/orders/{id}} & ID & Token válido & Devuelve pedido o \codefile{404}.\\
\bottomrule
\end{longtable}

\textbf{OrderCreate}
\begin{lstlisting}[language=json]
{
  "table_id": 1,
  "client_id": 2,
  "order_type": "local",
  "restaurant_id": 1
}
\end{lstlisting}

\textbf{Interpretación}

Un pedido representa la cabecera de la orden. Los ítems existen en el modelo y DAO, y son usados por exportaciones analíticas y de grafo, aunque no hay router CRUD público para \codefile{order\_items} en la API actual.

\subsection{Reservaciones}

\begin{longtable}{p{0.20\textwidth}p{0.22\textwidth}p{0.18\textwidth}p{0.32\textwidth}}
\toprule
\textbf{Endpoint} & \textbf{Entrada} & \textbf{Auth} & \textbf{Comportamiento}\\
\midrule
\endpoint{POST}{/reservations/} & \codefile{ReservationCreate} & Token válido & Crea reserva si mesa y cliente existen.\\
\endpoint{DELETE}{/reservations/{id}} & ID & Token válido & Cancela/elimina reserva. Responde \codefile{204}.\\
\bottomrule
\end{longtable}

\textbf{ReservationCreate}
\begin{lstlisting}[language=json]
{
  "table_id": 1,
  "client_id": 2,
  "reservation_date": "2026-06-25T19:30:00",
  "reservation_status": 1
}
\end{lstlisting}

\textbf{Interpretación}

\codefile{reservation\_date} se envía como fecha/hora ISO 8601. \codefile{reservation\_status} es entero no negativo y representa el estado de negocio definido por el sistema.

\subsection{Búsqueda}

El microservicio Search se publica normalmente bajo \codefile{/search}. Si se accede directo al contenedor, las rutas no incluyen prefijo.

\small
\begin{tabularx}{\textwidth}{
>{\raggedright\arraybackslash}p{0.26\textwidth}
>{\raggedright\arraybackslash}p{0.20\textwidth}
>{\raggedright\arraybackslash}p{0.17\textwidth}
>{\raggedright\arraybackslash}X
}
\toprule
\textbf{Endpoint} & \textbf{Entrada} & \textbf{Auth} & \textbf{Comportamiento}\\
\midrule
\endpoint{GET}{\codefile{/products?q=texto}} 
& Query \codefile{q} 
& Token válido 
& Busca productos en Elasticsearch por nombre, categoría o descripción. Usa Redis.\\

\endpoint{GET}{\codefile{/products/category/{category}}} 
& Categoría 
& Token válido 
& Filtra por categoría exacta normalizada. Usa Redis.\\

\endpoint{POST}{\codefile{/reindex}} 
& Ninguna 
& Token válido 
& Recrea índice \codefile{products}, indexa todos los menús y limpia caché de búsqueda.\\
\bottomrule
\end{tabularx}

\vspace{0.6cm}

\textbf{Respuesta de búsqueda}
\begin{lstlisting}[language=json]
[
  {
    "menu_id": 1,
    "dish_name": "Pizza Margarita",
    "category": "Plato fuerte",
    "description": "Pizza con queso, tomate y albahaca",
    "price": 7200.0,
    "restaurant_id": 1
  }
]
\end{lstlisting}

\textbf{Respuesta de reindexado}
\begin{lstlisting}[language=json]
{
  "message": "Productos reindexados correctamente",
  "total": 25
}
\end{lstlisting}

\subsection{Exportación para grafo}

\subsubsection{\endpoint{GET}{/graph/export}}

Exporta datos operacionales para Neo4J y para el generador de seed del Data Warehouse.

\textbf{Respuesta}
\begin{lstlisting}[language=json]
{
  "roles": [],
  "users": [],
  "restaurants": [],
  "products": [],
  "orders": [],
  "order_items": []
}
\end{lstlisting}

\textbf{Interpretación}: cada arreglo contiene registros normalizados. En \codefile{orders}, la API calcula \codefile{total} sumando \codefile{price * quantity} de sus ítems.

\section{Flujos de datos operacionales}

\subsection{Administrador agrega plato}

\begin{figure}[H]
    \centering
    \includegraphics[width=\textwidth]{documentation/Diagrama de flujo de datos - Admin_Menu.png}
    \caption{Flujo de datos para creación de menú por administrador.}
\end{figure}

El administrador envía una solicitud para agregar un platillo. Nginx redirige hacia FastAPI. La API valida el token con Keycloak y verifica rol \codefile{admin}. Después valida que el restaurante exista, crea el documento/fila en la base activa y limpia cachés afectadas. La respuesta JSON confirma el recurso creado.

\subsection{Cliente busca productos}

\begin{figure}[H]
    \centering
    \includegraphics[width=\textwidth]{documentation/Diagrama de flujo de datos - Cliente_Búsqueda.png}
    \caption{Flujo de datos para búsqueda de productos por cliente.}
\end{figure}

El cliente consulta productos. La API valida token, consulta Redis y, si hay \textit{cache hit}, responde inmediatamente. Si hay \textit{cache miss}, ejecuta búsqueda en Elasticsearch, procesa documentos encontrados, guarda el resultado en Redis y responde al cliente.

\subsection{Cliente crea reservación}

\begin{figure}[H]
    \centering
    \includegraphics[width=\textwidth]{documentation/Diagrama de flujo de datos - Cliente_Reserva.png}
    \caption{Flujo de datos para creación de reservación.}
\end{figure}

El cliente solicita una reserva. La API valida token, verifica que la mesa exista, verifica que el cliente exista y crea la reserva en la base activa. La respuesta confirma la reserva creada.

\section{Ejecución local}

\subsection{Levantar el sistema completo}

\begin{lstlisting}[language=bash]
docker compose up -d --build
\end{lstlisting}

\subsection{Verificar servicios}

\begin{lstlisting}[language=bash]
docker compose ps
docker compose logs -f keycloak
docker compose logs -f api
docker compose logs -f seed
\end{lstlisting}

\subsection{Credenciales por defecto}

\begin{longtable}{p{0.25\textwidth}p{0.25\textwidth}p{0.35\textwidth}}
\toprule
\textbf{Servicio} & \textbf{Usuario / password} & \textbf{URL}\\
\midrule
Keycloak Admin & \codefile{admin/admin} & \codefile{http://localhost:8001/admin/master/console/}\\
Usuario seed API & \codefile{admin/admin} & vía \codefile{/auth/login}\\
PostgreSQL & \codefile{postgres/postgres123} & \codefile{localhost:5432}\\
Superset & \codefile{admin/admin} & \codefile{http://localhost:8088}\\
Airflow & \codefile{admin/admin} & \codefile{http://localhost:8090}\\
Neo4J & \codefile{neo4j/restaurant123} & \codefile{http://localhost:7474}\\
\bottomrule
\end{longtable}

\section{Carpeta \codefile{dataW}: plataforma analítica completa}

\subsection{Objetivo de \codefile{dataW}}

\codefile{dataW} implementa la parte analítica del proyecto. Su propósito es transformar datos operacionales en información agregada para análisis OLAP y dashboards.

\begin{lstlisting}
PostgreSQL / MongoDB
        |
API /graph/export
        |
Generador HQL operacional
        |
Hive Data Warehouse
        |
Vistas OLAP
        |
Superset
        |
Dashboards
\end{lstlisting}

Airflow coordina el flujo completo y Spark ejecuta análisis distribuidos complementarios.

\subsection{Estructura de \codefile{dataW}}

\begin{longtable}{p{0.28\textwidth}p{0.62\textwidth}}
\toprule
\textbf{Ruta} & \textbf{Descripción}\\
\midrule
\codefile{dataW/warehouse/} & Esquema estrella, scripts Hive, seed operacional, ETL histórico y pruebas del Data Warehouse.\\
\codefile{dataW/spark/} & Job PySpark, salidas CSV/JSON y validación del procesamiento distribuido.\\
\codefile{dataW/visualization/} & Configuración de Superset, consultas SQL y catálogo de dashboards.\\
\codefile{dataW/airflow/} & DAG, scripts de orquestación, Dockerfile, estado del pipeline y pruebas de Airflow.\\
\bottomrule
\end{longtable}

\subsection{Data Warehouse en Hive}

El Data Warehouse usa un esquema estrella. Las dimensiones describen contexto; los hechos almacenan métricas.

\subsubsection{Dimensiones}

\begin{longtable}{p{0.22\textwidth}p{0.68\textwidth}}
\toprule
\textbf{Tabla} & \textbf{Propósito}\\
\midrule
\codefile{dim\_time} & Fechas, día, semana, mes, trimestre, año, fin de semana y temporada.\\
\codefile{dim\_customer} & Cliente, tipo, zona geográfica, nivel de lealtad, gasto total y actividad.\\
\codefile{dim\_product} & Producto, categoría, precio, costo, margen y disponibilidad.\\
\codefile{dim\_restaurant} & Restaurante, ubicación, zona, capacidad, contacto y estado.\\
\codefile{dim\_status} & Estados de órdenes y reservaciones.\\
\bottomrule
\end{longtable}

\subsubsection{Hechos}

\begin{longtable}{p{0.22\textwidth}p{0.68\textwidth}}
\toprule
\textbf{Tabla} & \textbf{Métricas}\\
\midrule
\codefile{fact\_orders} & Cantidad, precio unitario, total, descuento, neto, impuesto, final, hora de orden y hora de entrega.\\
\codefile{fact\_reservations} & Tamaño del grupo, duración, ocupación de mesa, no-show, fecha de reserva, check-in y check-out.\\
\bottomrule
\end{longtable}

\subsubsection{Vistas OLAP actuales}

El archivo \codefile{dataW/warehouse/schemas/hive\_dashboard\_views.sql} crea las vistas usadas por los dashboards de Superset:

\begin{itemize}
    \item \codefile{cubo\_ingresos\_mes\_categoria}
    \item \codefile{cubo\_actividad\_clientes\_zona}
    \item \codefile{cubo\_ordenes\_completadas\_canceladas}
\end{itemize}

\textbf{Ejemplo de consulta OLAP}
\begin{lstlisting}[language=SQL]
SELECT
    year,
    month_name,
    category,
    ingresos_totales
FROM cubo_ingresos_mes_categoria
ORDER BY year, month, category;
\end{lstlisting}

\subsection{Generación de seed operacional}

El servicio \codefile{warehouse-seed-generator} ejecuta:

\begin{lstlisting}
dataW/warehouse/scripts/generate_operational_seed.py
\end{lstlisting}

Este script consulta:

\begin{lstlisting}
http://api/graph/export
\end{lstlisting}

y genera:

\begin{lstlisting}
dataW/warehouse/generated/operational_seed.hql
\end{lstlisting}

El seed transforma usuarios, restaurantes, productos, pedidos e ítems en filas dimensionales y de hechos para Hive. Si faltan \codefile{order\_items}, el script puede completar datos derivados por restaurante para mantener un seed analítico utilizable.

\subsection{Spark}

El job \codefile{dataW/spark/jobs/analytics\_job.py} demuestra procesamiento distribuido con tres análisis:

\begin{enumerate}
    \item Tendencias de consumo por mes, categoría y producto.
    \item Horarios pico por día, hora y categoría.
    \item Crecimiento mensual con funciones ventana.
\end{enumerate}

\subsubsection{Entradas simuladas y escalamiento}

El job crea DataFrames base para tiempo, productos, restaurantes, órdenes y reservaciones. Luego escala hechos con:

\begin{lstlisting}
SPARK_SCALE_FACTOR=1000
\end{lstlisting}

Con el valor por defecto se generan 6000 órdenes y 5000 reservaciones.

\subsubsection{Salidas}

\begin{longtable}{p{0.30\textwidth}p{0.60\textwidth}}
\toprule
\textbf{Salida} & \textbf{Contenido}\\
\midrule
\codefile{output/consumption\_trends} & CSV de tendencias de consumo.\\
\codefile{output/peak\_hours} & CSV de horarios pico.\\
\codefile{output/monthly\_growth} & CSV de crecimiento mensual.\\
\codefile{output/validation\_summary.json} & Conteos de entrada y salida; evidencia de ejecución.\\
\bottomrule
\end{longtable}

\subsection{Superset}

Superset construye dashboards consultando directamente Hive. La herramienta no consume la API transaccional; consume vistas OLAP ya materializadas lógicamente.

\subsubsection{Dashboards requeridos}

\begin{longtable}{p{0.28\textwidth}p{0.30\textwidth}p{0.32\textwidth}}
\toprule
\textbf{Dashboard} & \textbf{Vista} & \textbf{Indicadores}\\
\midrule
Ingresos por mes y categoría & \codefile{cubo\_ingresos\_mes\_categoria} & Ingresos totales, órdenes, ingreso por categoría, serie mensual.\\
Actividad de clientes por zona & \codefile{cubo\_actividad\_clientes\_zona} & Clientes únicos, órdenes, ingresos por zona, ubicación.\\
Pedidos completados vs cancelados & \codefile{cubo\_ordenes\_completadas\_canceladas} & Cantidad por estado, tasa de cancelación, monto asociado.\\
\bottomrule
\end{longtable}

\subsection{Airflow}

Airflow ejecuta el DAG \codefile{restaurant\_olap\_pipeline}. Su schedule es diario:

\begin{lstlisting}
@daily
\end{lstlisting}

\subsubsection{Tareas del DAG}

\begin{lstlisting}
start
  -> extract_from_source
  -> run_spark_transformations
  -> load_to_data_warehouse
  -> validate_warehouse
  -> validate_delivery_routes_optional
  -> check_product_catalog_changes
  -> reindex_elasticsearch_if_needed / skip_reindex
  -> finish
\end{lstlisting}

\subsubsection{Descripción de tareas}

\begin{longtable}{p{0.30\textwidth}p{0.60\textwidth}}
\toprule
\textbf{Tarea} & \textbf{Función}\\
\midrule
\codefile{extract\_from\_source} & Verifica conectividad a PostgreSQL y MongoDB. Puede continuar si \codefile{ALLOW\_SOURCE\_UNAVAILABLE=true}.\\
\codefile{run\_spark\_transformations} & Ejecuta el job PySpark en un contenedor conectado al cluster Spark.\\
\codefile{load\_to\_data\_warehouse} & Crea base, esquema estrella, carga seed si corresponde y crea vistas OLAP.\\
\codefile{validate\_warehouse} & Verifica existencia y conteos positivos en hechos y cubos requeridos.\\
\codefile{validate\_delivery\_routes\_optional} & Ejecuta validación local de rutas si \codefile{ENABLE\_DELIVERY\_ROUTE\_VALIDATION=true}.\\
\codefile{check\_product\_catalog\_changes} & Calcula hash de \codefile{dim\_product} para saber si cambió el catálogo.\\
\codefile{reindex\_elasticsearch\_if\_needed} & Llama \codefile{POST http://search/reindex} si el catálogo cambió o si se fuerza.\\
\codefile{skip\_reindex} & Rama usada cuando no hay cambios.\\
\bottomrule
\end{longtable}

\section{Neo4J y asignación de rutas}

\subsection{Modelo de grafo}

El grafo contiene:

\begin{itemize}
    \item \codefile{Usuario}
    \item \codefile{Producto}
    \item \codefile{Pedido}
    \item \codefile{Restaurante}
    \item \codefile{Ubicacion}
    \item \codefile{Repartidor}
\end{itemize}

Relaciones principales:

\begin{itemize}
    \item \codefile{(:Usuario)-[:REALIZO]->(:Pedido)}
    \item \codefile{(:Pedido)-[:CONTIENE]->(:Producto)}
    \item \codefile{(:Pedido)-[:SALE\_DE]->(:Restaurante)}
    \item \codefile{(:Pedido)-[:ENTREGAR\_EN]->(:Ubicacion)}
    \item \codefile{(:Repartidor)-[:ASIGNADO\_A]->(:Pedido)}
\end{itemize}

\subsection{Heurística de vecino más cercano}

El algoritmo local:

\begin{enumerate}
    \item Carga repartidores, pedidos y ubicaciones desde JSON.
    \item Calcula distancia entre ubicaciones con Haversine.
    \item Procesa repartidores por rondas.
    \item Para cada repartidor elige el pedido pendiente más cercano.
    \item Actualiza la ubicación actual del repartidor.
    \item Continúa hasta agotar pedidos o capacidad.
\end{enumerate}

\section{Pruebas y validación}

\subsection{Pruebas unitarias}

\begin{lstlisting}[language=bash]
pytest backend/tests/unit
\end{lstlisting}

\subsection{Pruebas de integración}

\begin{lstlisting}[language=bash]
docker compose up -d --build
pytest backend/tests/integration
\end{lstlisting}

\subsection{Validación Spark}

\begin{lstlisting}[language=bash]
.\dataW\spark\tests\run_full_spark_validation.ps1
\end{lstlisting}

\subsection{Validación Data Warehouse}

\begin{lstlisting}[language=bash]
.\dataW\warehouse\tests\run_full_dw_validation.ps1
\end{lstlisting}

\subsection{Validación visualización}

\begin{lstlisting}[language=bash]
.\dataW\visualization\tests\run_full_visualization_validation.ps1
\end{lstlisting}

\subsection{Validación Airflow}

\begin{lstlisting}[language=bash]
python dataW\airflow\test\test_airflow_point4_requirements.py
\end{lstlisting}

\subsection{Validación Neo4J}

\begin{lstlisting}[language=bash]
python .\neo4j\delivery_assignment.py
python .\neo4j\test_delivery_routes.py
\end{lstlisting}

\section{Seguridad}

\begin{itemize}
    \item Las rutas protegidas requieren JWT válido.
    \item Las operaciones administrativas validan rol \codefile{admin}.
    \item La API no confía solo en datos del body; revisa existencia de dependencias en la base.
    \item Keycloak centraliza credenciales y emisión de tokens.
    \item Los tokens se validan criptográficamente contra JWKS.
    \item Nginx puede actuar como punto de control para exponer solo rutas necesarias.
\end{itemize}

\section{Consideraciones de consistencia}

\begin{itemize}
    \item El registro de usuario realiza dos escrituras: Keycloak y base local. Si falla la segunda, intenta compensar eliminando en Keycloak.
    \item El caché se invalida cuando se modifican menús porque afecta tanto listados como búsquedas.
    \item Elasticsearch no es la fuente de verdad; se reconstruye desde menús mediante \codefile{/reindex}.
    \item Hive tampoco es fuente transaccional; representa una copia analítica.
    \item Airflow evita duplicar cargas revisando si \codefile{fact\_orders} ya tiene datos.
\end{itemize}

\section{Limitaciones conocidas}

\begin{itemize}
    \item No existe router público completo para \codefile{order\_items}; el modelo y DAO sí existen.
    \item Las reservas no tienen endpoint \codefile{GET} ni \codefile{PUT} público en la implementación actual.
    \item Algunos cubos mencionados en guías históricas existen como diseño conceptual, pero el archivo activo de dashboards crea tres vistas OLAP requeridas.
    \item El proyecto de rutas usa ubicaciones proyectadas o simuladas porque la base operacional no almacena coordenadas reales de clientes/restaurantes.
    \item La validación de audience del JWT está desactivada porque Keycloak no incluye audience por defecto sin protocol mapper.
\end{itemize}

\section{Apéndice A: resumen de endpoints}

\small
\begin{longtable}{p{0.18\linewidth}p{0.22\linewidth}p{0.18\linewidth}p{0.34\linewidth}}
\toprule
\textbf{Método} & \textbf{Ruta} & \textbf{Entrada} & \textbf{Respuesta exitosa}\\
\midrule
GET & \codefile{/ping} & --- & \codefile{\{"message":"pong"\}}\\
GET & \codefile{/health} & --- & \codefile{\{"status":"ok"\}}\\
POST & \codefile{/auth/register} & \codefile{RegisterRequest} & \codefile{RegisterResponse}, 201\\
POST & \codefile{/auth/login} & \codefile{LoginRequest} & Tokens Keycloak\\
GET & \codefile{/users/me} & Token & \codefile{UserResponse}\\
PUT & \codefile{/users/{id}} & \codefile{UserUpdate} & \codefile{UserResponse}\\
DELETE & \codefile{/users/{id}} & ID & 204\\
GET & \codefile{/restaurants/} & Token & Lista \codefile{RestaurantResponse}\\
GET & \codefile{/restaurants/{id}} & ID & \codefile{RestaurantResponse}\\
POST & \codefile{/restaurants/} & \codefile{RestaurantCreate} & \codefile{RestaurantResponse}, 201\\
PUT & \codefile{/restaurants/{id}} & \codefile{RestaurantUpdate} & \codefile{RestaurantResponse}\\
DELETE & \codefile{/restaurants/{id}} & ID & 204\\
GET & \codefile{/menus/} & Token & Lista \codefile{MenuResponse}\\
GET & \codefile{/menus/{id}} & ID & \codefile{MenuResponse}\\
POST & \codefile{/menus/} & \codefile{MenuCreate} & \codefile{MenuResponse}, 201\\
PUT & \codefile{/menus/{id}} & \codefile{MenuUpdate} & \codefile{MenuResponse}\\
DELETE & \codefile{/menus/{id}} & ID & 204\\
GET & \codefile{/tables/} & Token & Lista \codefile{TableResponse}\\
GET & \codefile{/tables/{id}} & ID & \codefile{TableResponse}\\
POST & \codefile{/tables/} & \codefile{TableCreate} & \codefile{TableResponse}, 201\\
PUT & \codefile{/tables/{id}} & \codefile{TableUpdate} & \codefile{TableResponse}\\
DELETE & \codefile{/tables/{id}} & ID & 204\\
POST & \codefile{/orders/} & \codefile{OrderCreate} & \codefile{OrderResponse}, 201\\
GET & \codefile{/orders/{id}} & ID & \codefile{OrderResponse}\\
POST & \codefile{/reservations/} & \codefile{ReservationCreate} & \codefile{ReservationResponse}, 201\\
DELETE & \codefile{/reservations/{id}} & ID & 204\\
GET & \codefile{/products?q=} & Query \codefile{q} & Lista de productos indexados\\
GET & \codefile{/products/category/{category}} & Categoría & Lista de productos indexados\\
POST & \codefile{/reindex} & Token & Mensaje y total indexado\\
GET & \codefile{/graph/export} & --- & Exportación operacional para grafo/DW\\
\bottomrule
\end{longtable}

\section{Apéndice B: archivos fuente relevantes}

\small
\begin{tabularx}{\textwidth}{>{\raggedright\arraybackslash}X >{\raggedright\arraybackslash}X}
\toprule
\textbf{Archivo} & \textbf{Uso}\\
\midrule
\codefile{backend/app/api_main.py} & Aplicación FastAPI transaccional.\\
\codefile{backend/app/search_main.py} & Aplicación FastAPI de búsqueda.\\
\codefile{backend/database.py} & Selección de motor y dependencias de base.\\
\codefile{backend/dao.py} & DAO común para PostgreSQL y MongoDB.\\
\codefile{backend/app/routers/*.py} & Rutas REST.\\
\codefile{backend/schemas/*.py} & Contratos Pydantic.\\
\codefile{backend/models/*.py} & Modelos SQLAlchemy.\\
\codefile{dbs/mongo/*.js} & Inicialización, índices y sharding de MongoDB.\\
\codefile{dbs/postgres/init.sql} & Roles iniciales de PostgreSQL.\\
\codefile{dataW/warehouse/schemas/schema_star.sql} & Esquema estrella Hive.\\
\codefile{dataW/warehouse/schemas/hive_dashboard_views.sql} & Vistas OLAP activas.\\
\codefile{dataW/spark/jobs/analytics_job.py} & Procesamiento Spark.\\
\codefile{dataW/airflow/dags/restaurant_olap_pipeline.py} & DAG principal de Airflow.\\
\codefile{dataW/visualization/dashboards/dashboard_catalog.md} & Diseño de dashboards.\\
\codefile{neo4j/delivery_assignment.py} & Heurística local de rutas.\\
\bottomrule
\end{tabularx}

\section{Conclusión}

El proyecto integra un sistema transaccional robusto con una plataforma analítica completa. FastAPI centraliza reglas de negocio y contratos REST; Keycloak protege el acceso; PostgreSQL y MongoDB permiten persistencia relacional o documental; Redis y Elasticsearch optimizan lectura y búsqueda; Nginx simplifica exposición; y \codefile{dataW} convierte los datos operacionales en análisis mediante Hive, Spark, Superset y Airflow. La arquitectura permite demostrar bases distribuidas, APIs, caché, búsqueda, OLAP, orquestación y grafos dentro de un mismo ecosistema reproducible.

\end{document}

